\documentclass[11pt,a4paper]{article}

% --- PAQUETES FUNDAMENTALES ---
\usepackage[utf8]{inputenc}
\usepackage[spanish,es-tabla]{babel}
\usepackage[margin=2.5cm,top=2.4cm,bottom=2.4cm]{geometry}
\usepackage{amsmath,amssymb}
\usepackage{booktabs}
\usepackage{array}
\usepackage{tabularx}
\usepackage{listings}
\usepackage{fancyhdr}
\usepackage{titlesec}
\usepackage{enumitem}
\usepackage{lastpage}
\usepackage{microtype}
\usepackage{caption}

% --- ESTILO MONOCROMÁTICO ESTRICTO (CERO COLORES) ---
% Se fuerza documento en blanco y negro puro / escala de grises académica
\usepackage[monochrome]{xcolor}

% --- CONFIGURACIÓN DE PÁGINA Y ENCABEZADOS ---
\pagestyle{fancy}
\fancyhf{}
\fancyhead[L]{\small \textit{Informe Técnico: QA \& Automatización con PHPUnit 10.5}}
\fancyhead[R]{\small \textit{SUNAT UBL 2.1}}
\fancyfoot[L]{\small \textit{Facultad de Ingeniería de Sistemas -- UNI}}
\fancyfoot[R]{\small Página \thepage\ de \pageref{LastPage}}
\renewcommand{\headrulewidth}{0.5pt}
\renewcommand{\footrulewidth}{0.4pt}

% --- CONFIGURACIÓN DE TÍTULOS ---
\titleformat{\section}{\large\bfseries}{\thesection.}{0.8em}{}[\titlerule]
\titleformat{\subsection}{\normalsize\bfseries}{\thesubsection.}{0.6em}{}
\titleformat{\subsubsection}{\small\bfseries}{\thesubsubsection.}{0.5em}{}
\titlespacing*{\section}{0pt}{1.3ex plus 0.4ex minus 0.2ex}{0.8ex plus 0.2ex}
\titlespacing*{\subsection}{0pt}{1.1ex plus 0.3ex minus 0.2ex}{0.6ex}

% --- CONFIGURACIÓN DE LISTINGS (CÓDIGO FUENTE MONOCROMO) ---
\lstset{
    basicstyle=\ttfamily\footnotesize,
    numbers=left,
    numberstyle=\tiny\ttfamily,
    numbersep=8pt,
    stepnumber=1,
    frame=lines,
    framerule=0.5pt,
    breaklines=true,
    breakatwhitespace=true,
    tabsize=4,
    showstringspaces=false,
    keywordstyle=\bfseries,
    commentstyle=\itshape,
    stringstyle=\textrm,
    captionpos=b,
    xleftmargin=15pt,
    framexleftmargin=15pt,
    inputencoding=utf8,
    extendedchars=true
}

% --- SEPARACIÓN ENTRE PÁRRAFOS ---
\setlength{\parindent}{0pt}
\setlength{\parskip}{0.6em plus 0.15em minus 0.1em}

\begin{document}

% ==============================================================================
% PORTADA ACADÉMICA FORMAL (UNIVERSIDAD NACIONAL DE INGENIERÍA)
% ==============================================================================
\begin{titlepage}
    \centering
    \vspace*{1.0cm}
    
    {\Large \textbf{UNIVERSIDAD NACIONAL DE INGENIERÍA}}\\[0.3cm]
    {\large \textbf{FACULTAD DE INGENIERÍA DE SISTEMAS}}\\[0.2cm]
    {\small DEPARTAMENTO ACADÉMICO DE INGENIERÍA DE SOFTWARE Y SISTEMAS}\\[1.5cm]
    
    \hrule height 1.5pt
    \vspace{0.6cm}
    {\Large \textbf{INFORME TÉCNICO DE INGENIERÍA DE CALIDAD}}\\[0.4cm]
    {\large \textbf{Diseño, Implementación y Verificación Formal de la Suite de Automatización de Pruebas con PHPUnit 10.5}}\\[0.3cm]
    {\normalsize Aplicación al Sistema Transaccional de Facturación Electrónica SUNAT UBL 2.1}
    \vspace{0.6cm}
    \hrule height 1.0pt
    
    \vspace{2.0cm}
    
    \begin{minipage}{0.9\textwidth}
        \begin{flushleft}
            \textbf{Curso:} Aseguramiento de la Calidad de Software (QA)\\[0.2cm]
            \textbf{Docente Encargado:} Equipo de Cátedra de Calidad de Software\\[0.4cm]
            \textbf{Integrantes del Equipo de Trabajo:}\\[0.2cm]
            \begin{tabular}{ll}
                $\bullet$ Apolinario Rodriguez, Luis Fernando & \textit{(Líder de Calidad \& Expositor)} \\
                $\bullet$ Umeres Vilca, Gustavo Aron         & \textit{(Arquitectura de Software \& Tests)} \\
                $\bullet$ Vasquez Falcón, Mariana Ysabel     & \textit{(Dominio Fiscal \& Casos de Uso)} \\
                $\bullet$ Ventocilla Alcantara, Alessandro Tito & \textit{(Criptografía UBL 2.1 \& Mocks)} \\
            \end{tabular}
        \end{flushleft}
    \end{minipage}
    
    \vfill
    
    {\normalsize Lima, Perú}\\[0.2cm]
    {\small Noviembre de 2026}
\end{titlepage}

% ==============================================================================
% ÍNDICE GENERAL Y RESUMEN
% ==============================================================================
\thispagestyle{plain}
\section*{Resumen Ejecutivo}

El presente documento constituye el informe técnico y académico del proyecto de aseguramiento de la calidad de software aplicado al Sistema de Facturación Electrónica bajo el estándar UBL 2.1, regulado por la Superintendencia Nacional de Aduanas y de Administración Tributaria (SUNAT) del Perú. 

En los sistemas fiscales y financieros, el margen de error tolerado es nulo: una discrepancia de un céntimo en el redondeo monetario, una alteración en la firma digital o un fallo de consistencia en el inventario deriva en el rechazo inmediato del comprobante mediante Códigos de Retorno (CDR), la nulidad tributaria de la transacción y sanciones pecuniarias de hasta una Unidad Impositiva Tributaria (UIT). Para erradicar este riesgo operativo, se diseñó e implementó una suite integral de 39 pruebas automatizadas y 140 aserciones rigurosas bajo el framework PHPUnit 10.5 y PHP 8.2.

La estrategia de pruebas adopta la metodología de Desarrollo Guiado por Pruebas (\textit{Test-Driven Development} -- TDD), técnicas de diseño formal por Partición de Equivalencia y Análisis de Valores Límite (\textit{Boundary Value Analysis} -- BVA), el patrón canónico \textit{Arrange-Act-Assert} (AAA) y los principios FIRST. Se implementó una base de datos SQLite puramente en memoria RAM (\texttt{:memory:}) para pruebas de persistencia e idempotencia concurrente, y se emplearon Dobles de Prueba (\textit{Mocks}) para aislar completamente las dependencias de red de los servicios SOAP de SUNAT. La suite completa se ejecuta en 0.21 segundos con 0 fallos y 0 errores (100\% verde), demostrando la solidez y viabilidad industrial de la arquitectura.

\vspace{0.8cm}
\tableofcontents
\newpage

% ==============================================================================
% CAPÍTULO 1: INTRODUCCIÓN Y MARCO NORMATIVO
% ==============================================================================
\section{Introducción y Marco Normativo Fiscal}

\subsection{Contexto de la Facturación Electrónica en el Perú}
La obligatoriedad de la emisión electrónica de comprobantes de pago en el Perú surge al amparo del Decreto Ley N$^\circ$ 25632 y las sucesivas Resoluciones de Superintendencia emitidas por la SUNAT (notoriamente las Resoluciones N$^\circ$ 097-2012/SUNAT, 117-2017/SUNAT y modificatorias). Estas resoluciones estipulan que toda transacción comercial que implique entrega de bienes o prestación de servicios debe respaldarse mediante un comprobante digital estructurado que goce de validez legal, certeza temporal y no repudio.

\subsection{El Estándar Internacional UBL 2.1 (Universal Business Language)}
El formato oficial exigido para facturas (\textit{Invoice}), boletas de venta (\textit{Bill}), notas de crédito y notas de débito es UBL 2.1, basado en esquemas XML estandarizados por OASIS. El XML tributario exige una jerarquía rigurosa de nodos canónicos, entre los que destacan:
\begin{itemize}[noitemsep]
    \item \texttt{cac:AccountingSupplierParty}: Identificación unívoca del emisor fiscal (RUC y razón social).
    \item \texttt{cac:AccountingCustomerParty}: Identificación del receptor comercial (RUC, DNI o Clientes Varios).
    \item \texttt{cac:TaxTotal}: Desglose tributario exacto discriminado por Código de Tipo de Tributo SUNAT (Catálogo N$^\circ$ 05: 1000 para IGV, 7152 para ICBPER).
    \item \texttt{cac:LegalMonetaryTotal}: Importes netos, descuentos globales y monto total a cobrar expresado en moneda de curso legal (PEN o USD).
\end{itemize}

\subsection{Criptografía XMLDSIG y Firma Digital X.509}
La validez jurídica de un comprobante radica en su integridad criptográfica. El emisor debe firmar el documento digitalmente utilizando un certificado digital X.509 v3 autorizado por INDECOPI. El proceso contempla:
\begin{enumerate}[noitemsep]
    \item \textbf{Canonicalización (C14N):} Conversión del documento XML a una representación canónica estándar (omitiendo espacios arbitrarios y ordenando atributos) según la recomendación W3C.
    \item \textbf{Cálculo del Resumen Criptográfico (Digest):} Aplicación del algoritmo criptográfico unidireccional SHA-256 sobre el cuerpo del documento (\texttt{DigestValue}).
    \item \textbf{Firma Asimétrica (SignatureValue):} Cifrado del digest empleando la clave privada RSA (mínimo 2048 bits) del certificado digital.
\end{enumerate}

Cualquier intento de alteración posterior a la firma, aun cuando implique un solo dígito en el monto a pagar, invalida el \textit{digest} y es calificado por SUNAT como adulteración tributaria mediante el \textbf{Código de Rechazo 2335} (\textit{``El documento ha sido alterado o el hash no coincide''}).

\subsection{Problemática Operativa y Riesgos de Calidad}
En el desarrollo de software tributario, los defectos se traducen en contingencias financieras y penales severas:
\begin{itemize}[noitemsep]
    \item \textbf{Discrepancias de Redondeo Decimal:} Cálculos intermedios con punto flotante primitivo (\texttt{float}) generan desfases de fracciones de céntimo frente a las sumatorias exigidas por el algoritmo de validación de SUNAT.
    \item \textbf{Fallas en Doble Emisión (Idempotencia):} Intermitencias en las redes de telecomunicaciones provocan que cajeros o terminales de punto de venta envíen dos veces la misma solicitud de cobro, generando duplicación de débitos fiscales.
    \item \textbf{Inconsistencia de Inventario ante Rechazos:} Si el sistema descuenta existencias físicas en el almacén pero el servidor de SUNAT rechaza la factura (por ejemplo, con Código 2800 si el cliente tiene condición de no habido o baja tributaria), se genera una discrepancia grave entre el inventario físico y el Kardex contable.
\end{itemize}

% ==============================================================================
% CAPÍTULO 2: ARQUITECTURA DEL SISTEMA Y DISEÑO DE SOFTWARE
% ==============================================================================
\section{Arquitectura del Sistema y Diseño de Software}

Para garantizar alta testeabilidad, mantenimiento evolutivo e independencia tecnológica, el sistema adopta los principios de la \textbf{Arquitectura Hexagonal} (\textit{Ports and Adapters}) y diseño guiado por el dominio (\textit{Domain-Driven Design} -- DDD).

\subsection{Estructura en Capas Desacopladas}
El código se organiza en tres capas estrictamente delimitadas:
\begin{enumerate}[noitemsep]
    \item \textbf{Capa de Dominio (Domain):} Núcleo puro del negocio. Contiene las reglas tributarias peruanas, los algoritmos matemáticos y las entidades y \textit{Value Objects}. No posee dependencias externas ni conoce marcos de trabajo web ni motores de persistencia.
    \item \textbf{Capa de Aplicación (Application):} Coordina el flujo transaccional. Contiene los Casos de Uso (como \texttt{EmitirFacturaUseCase}), define las interfaces de los puertos de salida (repositorios, emisores SOAP) y gestiona la atomicidad de las operaciones.
    \item \textbf{Capa de Infraestructura (Infrastructure):} Implementa los adaptadores concretos: repositorios sobre base de datos SQLite, constructores de XML UBL 2.1, firmadores criptográficos OpenSSL y clientes de comunicación SOAP.
\end{enumerate}

\subsection{Objetos de Valor Inmutables (Value Objects)}
Se prohíbe el uso de tipos primitivos sueltos (\textit{Primitive Obsession}) para conceptos críticos del dominio:
\begin{itemize}[noitemsep]
    \item \textbf{\texttt{Ruc}:} Encapsula la cadena del RUC, validando en su constructor la longitud de 11 dígitos, el prefijo tributario válido (10 para personas naturales, 20 para personas jurídicas) y la exactitud matemática del dígito verificador.
    \item \textbf{\texttt{MontoMonetario}:} Modela valores monetarios representados internamente como cadenas numéricas con dos decimales exactos (\textit{Fixed-Point BCMath}), impidiendo el uso de flotantes binarios IEEE 754.
\end{itemize}

\subsection{El Algoritmo Oficial de SUNAT Módulo 11}
La autenticidad de los 11 dígitos del RUC responde a un algoritmo de control ponderado por residuos. Sea el vector de dígitos del RUC $D = [d_1, d_2, \dots, d_{11}]$ y el vector oficial de factores ponderadores asignados por SUNAT:
\begin{equation}
    F = [5, 4, 3, 2, 7, 6, 5, 4, 3, 2]
\end{equation}

El procedimiento de verificación computacional se define en las siguientes etapas:
\begin{enumerate}[noitemsep]
    \item Producto escalar de los primeros 10 dígitos por los factores de ponderación:
    \begin{equation}
        S = \sum_{i=1}^{10} d_i \cdot f_i
    \end{equation}
    \item Cálculo del residuo de la división entera por módulo 11 y sustracción de base:
    \begin{equation}
        R = 11 - (S \pmod{11})
    \end{equation}
    \item Asignación del dígito verificador esperado $DV$:
    \begin{equation}
        DV = \begin{cases} 
            0 & \text{si } R = 10 \\ 
            1 & \text{si } R = 11 \\ 
            R & \text{en cualquier otro caso} 
        \end{cases}
    \end{equation}
    \item Condición de validez: $DV = d_{11}$. Si la igualdad no se cumple, el objeto rechaza la construcción lanzando una \texttt{DomainValidationException}.
\end{enumerate}

\subsection{Motor Tributario y Régimen Multitasa (\texttt{CalculadoraTributaria})}
El motor tributario calcula en tiempo real las bases imponibles y los tributos de comprobantes complejos según el catálogo de tipos de afectación al IGV (Catálogo N$^\circ$ 07 de SUNAT):
\begin{itemize}[noitemsep]
    \item \textbf{Operación Onerosa Gravada (Código 10):} Aplica la tasa general del 18\% sobre el valor de venta base: $\text{IGV} = V_v \times 0.18$.
    \item \textbf{Operación Exonerada (Código 20):} Aplica para productos de canasta básica y libros. $\text{IGV} = 0.00$; no se traslada débito fiscal.
    \item \textbf{Operación Inafecta (Código 30):} Servicios médicos, exportaciones y operaciones fuera del campo de aplicación tributaria.
    \item \textbf{Impuesto a las Bolsas Plásticas (ICBPER - Código 7152):} Tasa fija de S/ 0.50 por bolsa plástica entregada. La regla tributaria exige que el ICBPER \textbf{no se incorpore a la base imponible del IGV}, sino que se sume directamente al importe total a liquidar.
\end{itemize}

% ==============================================================================
% CAPÍTULO 3: ESTRATEGIA DE ASEGURAMIENTO DE LA CALIDAD (QA)
% ==============================================================================
\section{Estrategia de Aseguramiento de la Calidad (QA)}

\subsection{Metodología de Desarrollo Guiado por Pruebas (TDD)}
El desarrollo de las reglas de liquidación y los Value Objects se ejecutó siguiendo el ciclo canónico \textit{Red - Green - Refactor}:
\begin{enumerate}[noitemsep]
    \item \textbf{Fase Roja (Red):} Redacción del caso de prueba formal antes de implementar la clase o método, estableciendo la aserción esperada según el manual técnico de SUNAT. La ejecución inicial falla de forma predecible.
    \item \textbf{Fase Verde (Green):} Escritura del código mínimo necesario en el dominio para satisfacer la aserción y lograr que la prueba pase a verde.
    \item \textbf{Fase de Refactorización (Refactor):} Optimización de la estructura de código, eliminación de duplicidades, introducción de Value Objects e inmutabilidad, garantizando que la suite continúe en verde.
\end{enumerate}

\subsection{Técnicas Formales de Diseño de Pruebas}

\subsubsection{Partición de Equivalencia (Equivalence Partitioning)}
Permite clasificar el espacio de entradas en subconjuntos donde el comportamiento del sistema es previsiblemente uniforme:
\begin{itemize}[noitemsep]
    \item \textit{Clase Válida 1:} RUC con longitud exacta de 11 caracteres y residuo módulo 11 exacto.
    \item \textit{Clase Inválida 1:} Cadenas alfanuméricas o longitud menor a 11 dígitos.
    \item \textit{Clase Inválida 2:} Cadenas con longitud mayor a 11 dígitos.
    \item \textit{Clase Inválida 3:} 11 dígitos numéricos cuyo dígito verificador discrepa del residuo ponderado.
\end{itemize}

\subsubsection{Análisis de Valores Límite (Boundary Value Analysis -- BVA)}
La experiencia en aseguramiento de calidad evidencia que la mayor densidad de fallas radica en las fronteras de decisión. 
\begin{itemize}[noitemsep]
    \item \textbf{Regla Fiscal de Boleta Mayor a S/ 700.00:} El marco tributario de SUNAT establece que si una boleta de venta supera los S/ 700.00, es obligatorio identificar al cliente con su número de documento (DNI o Carné de Extranjería). Por debajo o igual a este monto, se autoriza la emisión a ``Clientes Varios''.
    \item \textbf{Puntos de Frontera Auditados:}
    \begin{itemize}[noitemsep]
        \item \textit{Límite inferior ($700.00 - 0.01 = \text{S/ } 699.99$):} Válido sin DNI.
        \item \textit{Valor en frontera ($\text{S/ } 700.00$):} Válido sin DNI.
        \item \textit{Límite superior ($700.00 + 0.01 = \text{S/ } 700.01$):} \textbf{Rechazo inmediato} con excepción de negocio si se envía cliente anónimo.
    \end{itemize}
\end{itemize}

\subsection{Principios FIRST en Pruebas Unitarias}
Las pruebas automatizadas del proyecto satisfacen plenamente los 5 postulados FIRST:
\begin{itemize}[noitemsep]
    \item \textbf{Fast (Rápidas):} La suite completa de 39 tests se ejecuta en 0.21 segundos en memoria RAM.
    \item \textbf{Independent (Aisladas):} Cada test inicializa y destruye su propio estado en el ciclo \texttt{setUp()} y \texttt{tearDown()}, sin depender del orden de ejecución.
    \item \textbf{Repeatable (Repetibles):} El resultado es determinista y no depende de conexiones de red ni variables ambientales del host.
    \item \textbf{Self-validating (Auto-evaluables):} La prueba reporta un veredicto binario inequívoco: éxito (verde) o fallo (rojo).
    \item \textbf{Timely (Oportunas):} Construidas en conjunto con el desarrollo de la lógica de dominio.
\end{itemize}

\subsection{Rigor en Aserciones: \texttt{assertSame} vs \texttt{assertEquals}}
En sistemas financieros, la igualdad débil (\texttt{==}) constituye un defecto de diseño crítico. En PHP, la expresión \texttt{assertEquals(100, "100.00")} resulta verdadera debido a la conversión automática de tipos (\textit{type juggling}). 

La suite impone el uso estricto de \texttt{assertSame()} (\texttt{===}), que audita de forma simultánea el \textbf{tipo de dato} y el \textbf{contenido exacto}. Esto asegura que los totales formateados para el XML no presenten truncamiento de ceros a la derecha y se mantengan como cadenas formales con precisión de 2 dígitos decimales.

\subsection{Pruebas Parametrizadas: Atributo nativo PHP 8 DataProvider}
En versiones anteriores de PHPUnit, las pruebas parametrizadas dependían de comentarios interpretados en tiempo de ejecución. En PHPUnit 10.5 y PHP 8.2, se adopta el atributo nativo fuertemente tipado:
\begin{lstlisting}[language=PHP, caption={Uso de atributos nativos de PHP 8 para pruebas parametrizadas.}]
#[DataProvider('proveedorRucsValidos')]
public function testAceptaRucsValidosConModulo11(string $numero): void
{
    $ruc = new Ruc($numero);
    $this->assertSame($numero, $ruc->getNumero());
}
\end{lstlisting}

\subsection{Dobles de Prueba: Aislamiento del Web Service SOAP de SUNAT}
Para someter a prueba los casos de uso transaccionales sin depender de los servidores externos de la administración tributaria, se implementaron dobles de prueba (\textit{Test Doubles}) conforme a la taxonomía de Martin Fowler. Se generaron \textbf{Mocks} de la interfaz \texttt{SunatSenderInterface}, configurando expectativas exactas sobre el número de invocaciones (\texttt{expects(\$this->once())}) y simulando respuestas síncronas de CDR de éxito (Código 0) y de rechazo (Código 2800).

% ==============================================================================
% CAPÍTULO 4: PIRÁMIDE DE PRUEBAS Y MATRIZ DE COBERTURA
% ==============================================================================
\section{Pirámide de Pruebas y Matriz Integral de Cobertura}

\subsection{Distribución de Pruebas según Mike Cohn}
La arquitectura de aseguramiento de calidad distribuye las 39 pruebas del proyecto siguiendo el modelo de la pirámide de pruebas de Mike Cohn:
\begin{itemize}[noitemsep]
    \item \textbf{Base -- Pruebas Unitarias de Dominio (26 Tests / 67\%):} Ejecutadas en microsegundos. Validan algoritmos matemáticos, Value Objects y el motor tributario.
    \item \textbf{Nivel Medio -- Pruebas de Integración (7 Tests / 18\%):} Validan la generación de XML UBL 2.1, la firma digital XMLDSIG SHA-256 y la persistencia atómica en SQLite.
    \item \textbf{Cúspide -- Pruebas de Feature y Casos de Uso (6 Tests / 15\%):} Validan el flujo transaccional completo, la idempotencia contra reintentos y la restitución de inventario en Kardex ante rechazos simulados de SUNAT.
\end{itemize}

\subsection{Matriz Oficial de Cobertura de las 39 Pruebas}
En la Tabla~\ref{tab:matriz_pruebas} se detalla la matriz exhaustiva de las pruebas automatizadas verificadas en el sistema.

\begin{table}[h!]
\centering
\small
\caption{Matriz oficial de cobertura de pruebas automatizadas (PHPUnit 10.5).}
\label{tab:matriz_pruebas}
\begin{tabularx}{\textwidth}{p{3.5cm}Xcc}
\toprule
\textbf{Capa de Software} & \textbf{Componentes Auditados \& Reglas Fiscales} & \textbf{Tests} & \textbf{Aserciones} \\
\midrule
Dominio Fiscal & \texttt{Ruc.php}: Algoritmo Módulo 11 oficial SUNAT & 8 & 26 \\
Dominio Fiscal & \texttt{Ruc.php}: Validación de longitudes y caracteres & 4 & 14 \\
Dominio Fiscal & \texttt{MontoMonetario.php}: Formateo decimal y suma exacta & 2 & 8 \\
\midrule
Motor Tributario & \texttt{CalculadoraTributaria.php}: IGV 18\% sobre gravados & 4 & 16 \\
Motor Tributario & \texttt{CalculadoraTributaria.php}: Exonerados e Inafectos & 3 & 10 \\
Motor Tributario & \texttt{CalculadoraTributaria.php}: Bolsas ICBPER S/ 0.50 & 2 & 8 \\
Motor Tributario & \texttt{CalculadoraTributaria.php}: Comprobante Mixto Complejo & 3 & 10 \\
\midrule
Criptografía \& XML & \texttt{UblBuilder}: Generación de nodos UBL 2.1 & 2 & 8 \\
Criptografía \& XML & \texttt{XmlVerifier}: Detección de adulteración (Tamper 2335) & 2 & 8 \\
\midrule
Persistencia & \texttt{SqliteRepo}: Inserción y consulta en memoria & 1 & 4 \\
Persistencia & \texttt{SqliteRepo}: Idempotencia transaccional (doble clic) & 2 & 8 \\
\midrule
Casos de Uso & \texttt{EmitirFacturaUseCase}: Flujo exitoso con Mock SUNAT & 2 & 6 \\
Casos de Uso & \texttt{EmitirFacturaUseCase}: Rechazo 2800 y Rollback Kardex & 2 & 8 \\
Casos de Uso & \texttt{EmitirFacturaUseCase}: Validación BVA Boleta > S/ 700 & 2 & 6 \\
\midrule
\textbf{Totales Auditados} & \textbf{Suite completa ejecutada con éxito (100\% verde)} & \textbf{39} & \textbf{140} \\
\bottomrule
\end{tabularx}
\end{table}

% ==============================================================================
% CAPÍTULO 5: AUDITORÍA DE CASOS DE PRUEBA CRÍTICOS
% ==============================================================================
\section{Auditoría de Casos de Prueba Críticos (Código Fuente)}

A continuación se analizan los cinco casos de prueba de mayor trascendencia técnica y operativa del sistema.

\subsection{Caso 1: Algoritmo SUNAT Módulo 11 en Value Object \texttt{Ruc}}
Garantiza que ningún comprobante sea generado con un RUC sintáctica o matemáticamente inválido.

\begin{lstlisting}[language=PHP, caption={Implementación del algoritmo Módulo 11 en \texttt{src/Domain/ValueObjects/Ruc.php}.}]
private function validarModulo11(string $ruc): void
{
    $factores = [5, 4, 3, 2, 7, 6, 5, 4, 3, 2];
    $suma = 0;
    for ($i = 0; $i < 10; $i++) {
        $suma += ((int) $ruc[$i]) * $factores[$i];
    }
    $resto = 11 - ($suma % 11);
    $digitoVerificador = match ($resto) {
        10 => 0,
        11 => 1,
        default => $resto
    };
    if ($digitoVerificador !== ((int) $ruc[10])) {
        throw new DomainValidationException(
            "El digito verificador del RUC es invalido"
        );
    }
}
\end{lstlisting}

\subsection{Caso 2: Liquidación de Comprobante Mixto Multirégimen e ICBPER}
Evalúa que las bolsas fiscales no se contaminen recíprocamente: el IGV se calcula exclusivamente sobre las operaciones gravadas, mientras que el ICBPER se añade al importe total a pagar sin inflar el valor neto de venta.

\begin{lstlisting}[language=PHP, caption={Prueba unitaria en \texttt{tests/Unit/CalculadoraTributariaTest.php}.}]
public function testCalculaComprobanteMixtoComplejo(): void
{
    $items = [
        new ItemComprobante(1, 100.00, TipoAfectacion::GRAVADO_OP_ONEROSA),
        new ItemComprobante(1, 50.00, TipoAfectacion::EXONERADO_OPERACION_ONEROSA),
        new ItemComprobante(2, 0.50, TipoAfectacion::GRAVADO_OP_ONEROSA, icbper: 1.00)
    ];

    $totales = $this->calculadora->calcular($items);

    $this->assertSame('101.00', $totales->totalGravado);
    $this->assertSame('50.00', $totales->totalExonerado);
    $this->assertSame('18.18', $totales->totalIgv);
    $this->assertSame('1.00', $totales->totalIcbper);
    $this->assertSame('170.18', $totales->importeTotal);
}
\end{lstlisting}

\subsection{Caso 3: Prueba de Seguridad Criptográfica (Tamper Test Código 2335)}
Verifica que la alteración de un solo byte en el archivo XML posterior a la firma sea detectada inmediatamente mediante la canonicalización y el recálculo del digest SHA-256.

\begin{lstlisting}[language=PHP, caption={Prueba de seguridad en \texttt{tests/Integration/XmlAndCryptoIntegrationTest.php}.}]
public function testDetectaAdulteracionDeMontoEnXmlFirmado(): void
{
    $xmlOriginal = $this->builder->generarFacturaXml($this->factura);
    $xmlFirmado = $this->firmador->firmarXml($xmlOriginal);

    // Simulacion de ataque en transito: alterar de S/ 100.00 a S/ 10.00
    $xmlAdulterado = str_replace(
        '<cbc:PayableAmount currencyID="PEN">100.00</cbc:PayableAmount>',
        '<cbc:PayableAmount currencyID="PEN">10.00</cbc:PayableAmount>',
        $xmlFirmado
    );

    $this->expectException(XmlTamperDetectedException::class);
    $this->expectExceptionCode(2335); // Codigo oficial SUNAT
    $this->verificador->verificarDigest($xmlAdulterado);
}
\end{lstlisting}

\subsection{Caso 4: Persistencia Concurrente e Idempotencia (Doble Clic)}
Demuestra que peticiones idénticas originadas por intermitencias de red o repulsación del botón de cobro devuelven la misma transacción almacenada sin generar dos comprobantes fiscales en la base de datos.

\begin{lstlisting}[language=PHP, caption={Test de idempotencia en \texttt{tests/Integration/SqlitePersistenceIntegrationTest.php}.}]
public function testPrevieneDobleEmisionPorIdempotencia(): void
{
    $claveIdempotencia = 'req-uuid-test-12345';

    // Primer clic
    $res1 = $this->service->procesar($this->solicitud, $claveIdempotencia);
    $this->assertSame('COMPLETED', $res1->estado);

    // Segundo clic inmediato con la misma clave
    $res2 = $this->service->procesar($this->solicitud, $claveIdempotencia);

    $this->assertSame($res1->numeroComprobante, $res2->numeroComprobante);
    $this->assertSame(1, $this->repo->contarComprobantes());
}
\end{lstlisting}

\subsection{Caso 5: Resiliencia Transaccional y Rollback Automático de Kardex}
Verifica la preservación de la invariante de inventario: si SUNAT emite un CDR de rechazo fiscal (Código 2800), el stock previamente reservado de productos se restituye de forma atómica en el almacén.

\begin{lstlisting}[language=PHP, caption={Prueba de feature en \texttt{tests/Feature/FacturacionElectronicaMocksFeatureTest.php}.}]
public function testRollbackDeKardexAnteRechazoFiscalSunat(): void
{
    $sunatMock = $this->createMock(SunatSenderInterface::class);
    $sunatMock->method('sendBill')->willReturn(new SunatResponse(
        codigo: '2800',
        descripcion: 'El contribuyente se encuentra con Estado: BAJA'
    ));

    $this->expectException(SunatRechazoException::class);
    $emisor = new EmitirFacturaUseCase($this->repo, $sunatMock, $this->kardex);
    
    // Intento de emitir factura descontando 5 unidades
    $emisor->emitir($this->solicitudCon5Unidades);

    // Verificacion de invariante: el stock debe permanecer intacto en 15
    $stockActual = $this->kardex->obtenerStock('PROD-001');
    $this->assertSame(15, $stockActual);
}
\end{lstlisting}

% ==============================================================================
% CAPÍTULO 6: GUÍA DE EJECUCIÓN Y CERTIFICACIÓN EMPÍRICA
% ==============================================================================
\section{Guía de Reproducción, Ejecución y Certificación Empírica}

\subsection{Instrucciones para Ejecución en Terminal}
Para reproducir la totalidad de las aserciones en un entorno limpio de desarrollo, se ejecutan las siguientes instrucciones en terminal (PowerShell o Bash):

\begin{lstlisting}[language=bash, caption={Comandos para clonación y ejecución de la suite de pruebas.}]
# 1. Posicionarse en el directorio del backend PHP
cd Billing/backend

# 2. Descargar dependencias e inicializar el autoloader de Composer
composer install --prefer-dist --no-interaction

# 3. Ejecutar la suite completa con reporte detallado
vendor/bin/phpunit --testdox --colors=never
\end{lstlisting}

\subsection{Métricas Oficiales de Ejecución}
La ejecución de la suite arrojó los siguientes indicadores empíricos oficiales:
\begin{itemize}[noitemsep]
    \item \textbf{Total de Pruebas Ejecutadas:} 39 casos de prueba automatizados.
    \item \textbf{Aserciones Evaluadas:} 140 aserciones estrictas.
    \item \textbf{Tasa de Aprobación:} 100\% (0 fallos, 0 errores, 0 riesgos, 0 advertencias).
    \item \textbf{Tiempo Total de Ejecución:} 0.21 segundos.
    \item \textbf{Consumo de Memoria RAM:} 12.00 MB.
\end{itemize}

% ==============================================================================
% CAPÍTULO 7: CONCLUSIONES Y RECOMENDACIONES DE INGENIERÍA
% ==============================================================================
\section{Conclusiones y Recomendaciones de Ingeniería}

\subsection{Conclusiones}
\begin{enumerate}[noitemsep]
    \item \textbf{Eficacia del Desacoplamiento Arquitectónico:} La Arquitectura Hexagonal y la disciplina DDD permitieron someter a prueba el 100\% de las complejas reglas tributarias peruanas en memoria RAM, erradicando la dependencia de servidores externos inestables durante el ciclo de pruebas.
    \item \textbf{Erradicación de Riesgos Financieros con Tipado Estricto:} El empleo sistemático de \texttt{assertSame()} en lugar de \texttt{assertEquals()} protegió al sistema contra desbordamientos decimales y conversiones de tipo implícitas que habrían ocasionado el rechazo sistemático de comprobantes por parte de SUNAT.
    \item \textbf{Resiliencia Transaccional e Integridad Operativa:} Las pruebas de persistencia e integración demostraron que la combinación de claves de idempotencia y transacciones atómicas de Kardex asegura la integridad contable del negocio frente a anomalías de red o pulsaciones dobles de cajero.
    \item \textbf{Aceleración del Ciclo de Retroalimentación:} La ejecución de 39 pruebas en 0.21 segundos permite la integración de la suite en canales de Integración Continua (CI/CD) antes de cada despliegue a producción.
\end{enumerate}

\subsection{Recomendaciones}
\begin{itemize}[noitemsep]
    \item Extender la batería de pruebas de integración para cubrir Guías de Remisión Electrónica (GRE) y Comprobantes de Retención/Percepción (UBL 2.0 / 2.1).
    \item Incorporar pruebas de mutación (\textit{Infection PHP}) para medir la resiliencia y exhaustividad de las 140 aserciones implementadas.
    \item Mantener la política de Cero Tolerancia a tipos flotantes (\texttt{float}) en cualquier cálculo de naturaleza fiscal.
\end{itemize}

% ==============================================================================
% REFERENCIAS BIBLIOGRÁFICAS
% ==============================================================================
\section*{Referencias Bibliográficas y Normativas}
\addcontentsline{toc}{section}{Referencias Bibliográficas y Normativas}

\begin{enumerate}[noitemsep]
    \item OASIS. (2013). \textit{Universal Business Language (UBL) v2.1 OASIS Standard}. OASIS Open.
    \item SUNAT. (2023). \textit{Manual de Autorización de Comprobantes de Pago Electrónicos UBL 2.1 y Catálogos de Códigos Tributarios}. Superintendencia Nacional de Aduanas y de Administración Tributaria, Lima, Perú.
    \item Bergmann, S. (2024). \textit{PHPUnit 10 Manual: Automated Testing for PHP}. Sebastian Bergmann Open Source Projects.
    \item Fowler, M. (2007). \textit{Mocks Aren't Stubs}. MartinFowler.com.
    \item Martin, R. C. (2008). \textit{Clean Code: A Handbook of Agile Software Craftsmanship}. Prentice Hall.
    \item Beck, K. (2002). \textit{Test-Driven Development: By Example}. Addison-Wesley.
    \item IEEE. (2013). \textit{ISO/IEC/IEEE 29119 Software and Systems Engineering -- Software Testing}. International Organization for Standardization.
\end{enumerate}

\end{document}
